\section{Related Work}
\label{sec:related_work}

\subsection{Static Model Merging}
Model merging has emerged as a key paradigm for fusing specialized, pre-trained neural network weights into generalist multi-task models without joint fine-tuning. Early foundational works, such as Model Soups~\cite{wortsman2022model} and Fisher Weighted Averaging~\cite{matena2021merging}, demonstrated that averaging the weights of models fine-tuned on the same dataset improves out-of-distribution performance. This was subsequently generalized to multi-task scenarios by Task Arithmetic~\cite{ilharco2022editing}, which computes task-specific vectors as the difference between expert and pre-trained weights and adds their linear combination to the base pre-trained model. To resolve parameter conflicts (such as sign disagreement and parameter redundancy), TIES-Merging~\cite{yadav2023ties} filters out small-magnitude task vector elements and resolves sign conflicts via consensus voting, while Dataless Knowledge Fusion~\cite{jin2022dataless} and SALSA~\cite{yu2023language} prune unimportant weights. Other approaches seek to align permutation symmetries of different experts before merging, such as Git Re-Basin~\cite{ainsworth2022git} and ZipIt!~\cite{stoica2023zip}. However, all these static methods rely on a manual, global scaling parameter and operate under the rigid, cold assumption of a flat Euclidean parameter landscape~\cite{ortizjimenez2023task}. This zero-temperature assumption causes severe performance drops when merging experts fine-tuned on highly diverse or heterogeneous tasks, as the straight-line average path is forced to cross high-loss, non-convex ridges that trigger representation collapse.

\subsection{Test-Time Adaptive Model Merging}
To overcome the limitations of manually tuned global scaling coefficients, test-time adaptive merging methods have been proposed. Under this setting, the merging coefficients are treated as trainable parameters and optimized directly on unlabeled target data streams during inference. AdaMerging~\cite{yang2024adamerging} first introduced this paradigm, optimizing layer-wise or task-wise coefficients by minimizing the entropy of the merged model's predictions on unlabeled calibration batches, drawing inspiration from standard unsupervised test-time adaptation methods like Tent~\cite{wang2020tent} and parameter-free TTA~\cite{boudiaf2022parameter}, as well as continual or single-sample test-time adaptation methods like CoTTA~\cite{wang2022continual}, MEMO~\cite{zhang2022memo}, and NOTE~\cite{gong2022note}. Recently, SyMerge~\cite{jung2025symerge} advanced this by using a self-labeling teacher alignment objective to stabilize training. More recently, FoldMerge~\cite{foldmerge_origami} proposed a non-linear coordinate-warping framework using RealNVP normalizing flows to map task vectors into a latent ``Origami Space'' where linear interpolation is safer.

However, as exposed in the \emph{Overfitting-Optimizer Paradox}~\cite{overfittingoptimizerparadox}, unsupervised test-time adaptive merging methods are plagued by a severe optimization pathology. Because calibration streams are small and fast-to-run, standard, unregularized entropy minimization (such as AdaMerging's objective) drives the model parameters into degenerate local minima with low entropy but zero generalizing representations. This causes catastrophic transductive overfitting, leading to complete representation collapse on complex domains. In contrast, ThermoMerge introduces a physically grounded, thermodynamic adaptation framework where the Helmholtz Free Energy serves as a natural, mathematically rigorous anchor that completely prevents transductive collapse.

\subsection{Deep Learning and Statistical Physics}
The cross-pollination between statistical mechanics and deep learning has a rich and storied history~\cite{carleo2019machine,bahri2020statistical}. Pioneering works modeled neural network loss surfaces as spin glasses, demonstrating that high-dimensional non-convex landscapes are filled with local minima and saddle points~\cite{mezard1987spin,choromanska2015loss}. Over the years, physical principles have been utilized across diverse deep learning domains, ranging from modern Hopfield networks connecting associative memories to self-attention~\cite{ramsauer2020hopfield}, to the Information Bottleneck principle framing optimal representation compression~\cite{tishby2015deep}. To navigate these rugged landscapes, researchers have drawn on physical processes: simulated annealing uses physical cooling to find global minima~\cite{kirkpatrick1983optimization}, while stochastic gradient descent (SGD) has been formally modeled as approximate Langevin diffusion or Bayesian inference~\cite{mandt2017stochastic}. Furthermore, Entropy-SGD~\cite{chaudhari2019entropy}, Hochreiter's early flat minima theory~\cite{hochreiter1997flat}, and physical theories of subdominant dense basins~\cite{baldassi2015subdominant} demonstrated that biasing optimization toward wide, flat valleys improves generalizability. 

While these works apply statistical physics to standard network training, ThermoMerge is, to the best of our knowledge, the first framework to translate physical thermodynamics directly to the challenge of model merging. By mapping classification logits to Boltzmann states and optimizing a physical Free Energy discrepancy under a simulated thermal schedule, we establish a robust, physically principled bridge that allows model merging to navigate highly non-convex multi-task loss frontiers.
